\clearpage
\fancyhead[L]{\small CONNAISSANCE DES TEMS FOR 1818.}
\hypertarget{native-600-lower}{}
\newcommand{\ObservedJupiterErrorMark}{\hbox{\vrule width0pt height4.000000bp depth0pt\special{pdf:bcontent}\special{pdf:code q 0 g 0.500000 4.000000 m 0.000000 3.500000 l 0.000000 3.250000 l 0.750000 2.500000 l 0.750000 2.250000 l 1.000000 2.000000 l 1.000000 1.750000 l 0.750000 1.500000 l 0.250000 1.500000 l 0.000000 1.250000 l 0.000000 0.500000 l 0.250000 0.250000 l 1.250000 0.250000 l 1.500000 0.000000 l 2.000000 0.000000 l 2.250000 0.250000 l 2.500000 0.250000 l 3.250000 1.000000 l 3.250000 1.250000 l 2.500000 2.000000 l 2.500000 2.500000 l 3.000000 3.000000 l 3.000000 3.250000 l 2.250000 4.000000 l 1.750000 4.000000 l 1.500000 3.750000 l 1.250000 3.750000 l 1.000000 4.000000 l h f* Q}\special{pdf:econtent}\kern3.500000bp}}

\begin{center}\small Göttingische gelehrte Anzeigen. No. 54. Pages 531..536. 5 April 1817.\\[3pt]\rule{.72\linewidth}{.3pt}\end{center}

\emph{Connaissance des tems, ou des mouvemens célestes, à l'usage des astronomes et des navigateurs, pour l'an \textup{1818}; publiée par le bureau des longitudes.} Paris 1815. Published by the widow Courcier. 412 pages in octavo.

The arrangement of the astronomical calendar has undergone a few small changes in this volume. The Moon's right ascension, formerly given only in arcminutes, is now calculated to seconds for noon and midnight. The declinations, by contrast, were formerly given every six hours and are now given in minutes for noon and midnight only. This change must be welcome to those astronomers who wish to determine the difference in their geographical longitudes through corresponding observations of lunar culminations. Furthermore, the individual phenomena, formerly listed separately for each month, are now more conveniently collected at the end of the calendar. The space thus gained has been used to give the heliocentric longitudes and latitudes of the planets and their right ascensions in units of time. It is only to be regretted that the right ascensions have been tabulated only to the minute (that is, to quarter degrees of arc). In order to judge in advance the positions of the planets relative to neighbouring fixed stars, it is very often desirable to know the positions of the planets in advance to the arcminute. It would be far more important to give the planets' distance from the Earth than to give the heliocentric positions, which are of little or no use in the ephemeris. It is much to be regretted that not a single one of our astronomical ephemerides gives this element, even though the distance from the Earth is equally important for aberration, parallax, and the reduction of measured diameters.

The accompanying papers begin with a number of extracts that Mr.\ \textsc{Delambre} has made from various English writings. Since these still reach few readers on the Continent, many will welcome the extracts. The accompanying remarks of so experienced a practitioner are also a pleasure to read, even if he often repeats himself. The following paper, by contrast, is written at tiresome and needless length. It concerns a problem of \textsc{Regiomontanus}, to state the circumstances under which two points on the ecliptic differ by as much in longitude as in right ascension. This slight and \SourceLeaf{601}easy problem did not deserve 10 pages in the Connoissance des tems. --- New tables for the aberration of the planets in longitude and latitude, by \textsc{Puissant}. The arrangement of these tables is as convenient as the nature of the subject allows, and they may be useful when the planet’s distance from the Earth is entirely unknown. The author remarks on the formulas used to compute the tables, Ces formules sont \emph{complètes}; il est remarquable en outre, qu'elles sont exactes, \emph{aux quantités près du troisième ordre}; car les termes en $\mathrm{e}^2$ se detruisent mutuellement, comme je m'en suis assuré. This rather obscure passage presumably means that the author has expanded the aberration in powers of the eccentricity, and that both the terms independent of the eccentricity and the first-order terms in the eccentricity are complete. If this interpretation is correct (without it the underlined expressions appear contradictory), the author seems to have overlooked that his formulas are in fact accurate beyond terms of third order in the eccentricity and possess \emph{absolute} completeness. --- By the same author, a note on a paper by \textsc{Lagrange} concerning the calculation of parallax in eclipses. Mr.\ \textsc{Puissant} uses auxiliary means in an attempt to put the formulas given by \textsc{Lagrange} for the apparent distance between the centres of the Sun and the Moon into a form better suited to numerical calculation. It is more convenient here, however, to use the table that gives at once the logarithms of sums and differences of quantities given directly by their logarithms. --- The new general table for the parabolic motion of comets by \textsc{Burckhardt} deserves preference over the ordinary table and even over \textsc{Barker}’s. It differs from the ordinary table in taking as its argument the logarithm of the time elapsed since passage through perihelion in place of the time itself. Compared with the ordinary table it saves a logarithm lookup. It has the same advantage over \textsc{Barker}’s table up to $45^\circ$ of true anomaly, and interpolation in \textsc{Burckhardt}’s table is also usually somewhat more convenient. --- New determination of Vesta’s orbit by \textsc{Daussy}. Mr.\ \textsc{Daussy} has rendered a great service by calculating the perturbations of Vesta’s motion due to Jupiter, Saturn and Mars using \textsc{Laplace}’s method, which may be considered adequate for that planet’s moderate eccentricity and inclination. The first seven oppositions do indeed show very satisfactory agreement. --- On a new apparatus for comparing linear measures, by \textsc{Prony}. The chief component of this apparatus is a millimetre divided on glass into 100 parts. Several French instrument makers produce this graduation with truly admirable fineness and precision. The measures to be compared must be combined in such a way that they differ by no more than a couple of millimetres. The apparatus determines only this difference. A fuller description would take us too far here. --- On the ebb and flow of the sea, by \textsc{Laplace}. This short paper is the introduction to new and interesting investigations by this great geometer into the observations of the tides made at Brest. It must make us very eager for the details of his investigations. The determinations derived from them for the mass of the Moon $(= \frac{1}{68.7})$ and the constant of nutation $(= 9''65)$ are particularly remarkable, all the more so because the observations of the pole star discussed so carefully by Mr.\ \textsc{von Lindenau} yielded a considerably smaller value of the latter. --- The next two papers by the same author, on the application of probability theory to the natural sciences, are concerned with the theory of the accuracy of results obtained by the method of least squares and with determining the accuracy of the observations themselves. The results of this investigation agree in substance entirely with what \textsc{Gauss} developed in Theoria Motus \SourceLeaf{602}Corporum Coelestium and in a paper in the Zeitschrift für Astronomie, although \textsc{Laplace} considered the subject from a somewhat different point of view. Several interesting applications of this theory are also given here. One is the determination of the mass of Jupiter from the perturbations that Saturn undergoes under Jupiter’s influence. The result \textsc{Bouvard} derived from this source is $\frac{1}{1070.5}$, and the probable error of this determination is found to be $= \pm \frac{1}{72\ObservedJupiterErrorMark}$ of that value. In other words, the true value is just as likely to lie between $\frac{1}{1069}$ and $\frac{1}{1072}$ as it is to fall outside those limits. \textsc{Laplace} preferred to give the extraordinarily small probability of exceeding somewhat wider limits. One can bet 1 against nearly a million that the error amounts to no more than $\frac{1}{100}$ of the whole. The most remarkable circumstance now arises. The mass of Jupiter derived by \textsc{Gauss} from his theory of the perturbations of Pallas by Jupiter differs considerably from that determination and falls well outside those limits. The main point is that probability theory, when applied to this determination, assigns it very nearly as small an uncertainty. This last point should come as no surprise, because the much stronger influence of Jupiter already makes up for the still smaller number of observations of the motions of Pallas. How are we to reconcile this contradiction, and what is the true mass of Jupiter? Continued observations of Pallas and an investigation of the perturbations of Juno will provide further information in the future. It seems reasonable, however, to doubt whether \textsc{Laplace}’s method of calculating perturbations, applied to Saturn, yields results quite as accurate as such a delicate investigation requires. --- On \textsc{Nonius}’s formulas for twilight, by \textsc{Delambre}. \textsc{Nonius}’s contribution to the problems concerning the duration of twilight has so far received too little recognition. In this very detailed paper Mr.\ \textsc{Delambre} shows that \textsc{Nonius} did in fact solve all the relevant problems very well, some of them better than his successors. We should regard this paper as a sample of a large work that Mr.\ \textsc{Delambre} intends to publish as a history of astronomy. It is to offer extracts from and commentary on all known books, at least all those that contain anything useful or noteworthy. Some are to be reprinted and all are to be presented at least in a new form. --- The last paper in this volume contains the elements and perturbation equations used to calculate \textsc{Bouvard}’s new tables for Saturn. It also compares the tables with all observed oppositions and quadratures of the planet since 1747, 130 in all. We should very much have liked all the observed positions underlying the tables to have been printed with them. However well the tables represent the observations, sooner or later further improvements will have to be sought, and for that purpose it will be necessary to go back to all those oppositions.
